\reportheading{4.\quad 任务 2：排烟设计与优化}
\subhead{4.1\quad 原布局的 FDS 筛选}
在 OP-6、基准烟炱产率下，16 个原布局方案中净高余量最大的组合为 0.16~m 挡烟垂壁和 0.32~m$^3$/s 排烟量。该组合最低净高为 0.095~m，对应余量为 $-0.285$~m（阈值 0.38~m，附录 C）。根据这组 FDS 结果，下一轮设计调整了补风路径和烟气捕集位置。

\subhead{4.2\quad 定位越限位置与调整补风口}
一组 OP-6 试验采用 0.36~m 高端墙补风口、0.24~m 挡烟垂壁及 1.28~m$^3$/s 总排烟量。在烟炱产率 0.05 下，远端测点（$x=6.1$~m）在 83.0~s 时降至 0.353~m，相邻前后两个样本均高于 0.60~m。按 0.38~m 的热分层判据，这一采样时刻越限。同期半高处能见度为 30~m，局部竖向温度剖面分层较弱。

在原网格上固定 1.28~m$^3$/s，在 OP-6 的两个烟炱产率下直接比较 0.44 与 0.48~m 高的端墙补风口。增高补风口后，产率 0.02 时的最低界面由 0.431 升至 0.583~m，产率 0.05 时则由 0.524 降至 0.497~m。两工况中较不利的最低高度由 0.431 提高到 0.497~m，因此后续对 OP-6 采用 0.48~m。六工况复核时，OP-1 采用 0.44~m、OP-6 采用 0.48~m 补风口，垂壁下垂深度统一为 0.24~m。
\reportfigure{\begin{tikzpicture}
\begin{axis}[width=.94\linewidth,height=5.6cm,xmin=.009,xmax=.112,ymin=.22,ymax=.405,
xtick={.02,.05,.10},ytick={.24,.28,.32,.36,.38,.40},scaled ticks=false,
xticklabel style={/pgf/number format/fixed,/pgf/number format/fixed zerofill,/pgf/number format/precision=2},
yticklabel style={/pgf/number format/fixed},axis lines*=left,ymajorgrids,grid style={ink!13},
xlabel={烟炱产率 (kg/kg)},ylabel={最低高度 (m)},legend columns=3,legend style={at={(.5,1.03)},anchor=south,draw=none,/tikz/every even column/.append style={column sep=6pt}}]
\addplot[only marks,teal,mark=*,mark size=3pt] coordinates {(0.02,0.2627491) (0.05,0.37114426) (0.1,0.24450079)};\addlegendentry{0.16 m 垂壁}
\addplot[only marks,rust,mark=triangle*,mark size=3pt] coordinates {(0.02,0.36496348) (0.05,0.3527722) (0.1,0.36869806)};\addlegendentry{0.24 m 垂壁}
\addplot[only marks,ink,mark=square*,mark size=3pt] coordinates {(0.10,0.26696128)};\addlegendentry{0.32 m 垂壁}
\addplot[ink!75,dashed,forget plot] coordinates {(.009,.38) (.112,.38)};
\node[anchor=south west,font=\sffamily\footnotesize,text=ink] at (axis cs:.011,.38){0.38 m 阈值};
\end{axis}\end{tikzpicture}}{图 3.\quad 1.28~m$^3$/s 下七个已完成 OP-6 垂壁工况的最低热分层高度。每个标记取八个测点在 120~s 内的最低值。0.32~m 垂壁仅在烟炱产率 0.10~kg/kg 下测试。}

\clearpage
\subhead{4.3\quad 设计风量的选取}
在双顶棚排烟口、0.24~m 挡烟垂壁及随门态预设的端墙补风口布局下，取初始两档风量（0.96 和 1.28~m$^3$/s）的中点 1.12~m$^3$/s 进行计算。设计需要回答的是，该风量能否在两种门态和三种烟炱产率下，将热分层界面维持在房间半高以上。
\reportfigure{\begin{tikzpicture}
\pgfplotsset{flowpanel/.style={
 scale only axis,height=3.35cm,ymin=.35,ymax=.62,ytick={.38,.44,.50,.56,.62},
 scaled ticks=false,yticklabel style={/pgf/number format/fixed,/pgf/number format/precision=2},
 ymajorgrids,grid style={ink!13},tick align=outside,title style={font=\sffamily\small\bfseries},
 tick label style={font=\sffamily\scriptsize},label style={font=\sffamily\scriptsize}}}
\begin{axis}[flowpanel,name=original,width=.47\linewidth,xmin=.55,xmax=6.45,
 title={(a) 原网格：六工况筛选},ylabel={最低净高 (m)},
 xtick={1,2,3,4,5,6},xticklabels={0.02,0.05,0.10,0.02,0.05,0.10},
 clip=false,legend style={at={(.50,.99)},anchor=north,font=\sffamily\scriptsize,fill=white},
 legend columns=2]
\addplot[only marks,mark=*,teal,mark size=1.9pt] coordinates {(1,.45742963) (2,.47169434) (3,.44104958) (4,.39573017) (5,.52617570) (6,.40609647)};
\addlegendentry{0.96}
\addplot[only marks,mark=triangle*,rust,mark size=2.4pt] coordinates {(1,.44252295) (2,.50930185) (3,.50127480) (4,.58267633) (5,.49747471) (6,.55033156)};
\addlegendentry{1.28}
\addplot[ink!70,dashed] coordinates {(.55,.38)(6.45,.38)};
\node[anchor=north,font=\sffamily\scriptsize] at (axis description cs:.25,-.15) {OP-1};
\node[anchor=north,font=\sffamily\scriptsize] at (axis description cs:.75,-.15) {OP-6};
\end{axis}
\begin{axis}[flowpanel,at={(original.east)},anchor=west,xshift=.065\linewidth,
 width=.34\linewidth,xmin=.65,xmax=3.35,title={(b) 细网格：OP-6，产率 0.02},
 xtick={1,2,3},xticklabels={0.96,1.12,1.28},xlabel={排烟量 (m$^3$/s)},yticklabels=\empty]
\addplot[ink!70,dashed] coordinates {(.65,.38)(3.35,.38)};
\addplot[only marks,mark=*,teal,mark size=2.3pt] coordinates {(1,.37169)};
\addplot[only marks,mark=diamond*,ink,mark size=3pt] coordinates {(2,.5522)};
\addplot[only marks,mark=triangle*,rust,mark size=2.6pt] coordinates {(3,.5429)};
\node[anchor=south,font=\sffamily\scriptsize] at (axis cs:2,.558) {0.5522 m};
\end{axis}
\end{tikzpicture}}{图 4.\quad 两阶段 FDS 风量比较，虚线为 0.38~m 限值。(a) 原网格比较 0.96 和 1.28~m$^3$/s，横轴按门态和烟炱产率（kg/kg）分组。(b) 在同一细网格和工况下，0.96~m$^3$/s 在 98~s 越限即停止；1.12 和 1.28~m$^3$/s 均完成 120~s。1.12~m$^3$/s 的六工况复核见图 5。}


原网格上，两档参照风量的最小余量都出现在基准烟炱产率（图 4a）。0.96~m$^3$/s 的余量为 15.7~mm，小于 40~mm 的竖向单元高度。将单元高度减半后，OP-6 基准烟炱工况的远端测柱在 98.0~s 记录到 0.3717~m，筛选在该时刻结束（表 D2）。图 4(b) 的三个点采用相同竖向分辨率，可直接比较。

\subhead{1.12~m$^3$/s 下的净高与功率}
在 OP-6、基准烟炱产率下，1.12~m$^3$/s 工况的 120~s 最低热分层界面为 0.5522~m，1.28~m$^3$/s 时为 0.5429~m。在拟选风量下，每条短支管承担 0.56~m$^3$/s，对应风速 1.75~m/s、估算压降 3.70~Pa。按假定效率 0.60，两条短支管合计估算电功率为 6.91~W，比 1.28~m$^3$/s 时的 10.32~W 低约 33\%。据此将 1.12~m$^3$/s 带入完整的六工况复核。

\clearpage
\subhead{六工况 FDS 复核}
在 1.12~m$^3$/s 下，六组竖向加密的计算均完成 120~s。火源名义功率为 106.2~kW，各组在 90--120~s 的实际平均热释放率介于 102.1 至 103.9~kW。八个中心面测柱记录的热分层界面均高于 0.38~m（表 D3）。最小余量为 55.8~mm，出现在 OP-6、烟炱产率 0.05。远端测柱（$x=6.1$~m）在 113.0~s 记录到 0.4358~m，同处纵向速度在 $z=0.42$ 与 0.44~m 之间变号。图 A2 给出同一时刻半高处的局部烟气分布。

% Rows are the six refined-grid FDS station minima at 1.12 m3/s.
% PINN points are the forecasts frozen in hybrid_frozen_1p12.json.
\reportfigure{\begin{tikzpicture}
\begin{axis}[width=.96\linewidth,height=5.1cm,xmin=.55,xmax=6.45,
 ymin=.355,ymax=.61,xtick={1,2,3,4,5,6},
 xticklabels={{OP-1\\0.02},{OP-1\\0.05},{OP-1\\0.10},{OP-6\\0.02},{OP-6\\0.05},{OP-6\\0.10}},
 xticklabel style={align=center},ytick={.38,.44,.50,.56,.60},ylabel={最低净高 (m)},
 ymajorgrids,unbounded coords=discard,
 legend style={at={(.5,1.03)},anchor=south,legend columns=3,draw=none},
 /tikz/every even column/.append style={column sep=5pt}]
\addplot[only marks,mark=*,teal,mark size=2.7pt]
 table[col sep=comma,x=position,y=fds_m]{data/verified-flow-forecast.csv};
\addlegendentry{细网格 FDS}
\addplot[only marks,mark=diamond*,rust,mark options={fill=white},mark size=3.2pt]
 table[col sep=comma,x expr=\thisrow{position}+.10,y=pinn_m]{data/verified-flow-forecast.csv};
\addlegendentry{冻结 PINN}
\addplot[ink,densely dashed] coordinates {(.55,.38)(6.45,.38)};
\addlegendentry{净高阈值}
\draw[teal,thin] (axis cs:5,.385)--(axis cs:5,.430);
\node[anchor=west,font=\sffamily\scriptsize,text=teal] at (axis cs:5.04,.408) {55.8 mm};
\end{axis}\end{tikzpicture}}{图 5.\quad 1.12~m$^3$/s 的六工况复核。横轴为门态和烟炱产率（kg/kg），实心点为八个测点在 120~s 内的 FDS 最低净高。空心菱形仅标出运行前已冻结的三组 PINN 预测。}


\subhead{4.4\quad CO 生成与局部烟气条件}
按指定产率，每消耗 1~kg 燃料生成 10~g CO。在早期 OP-6 试验中，排烟量为 1.28~m$^3$/s，垂壁深度为 0.24~m。根据热释放率积分和模型燃烧热，随烟炱产率升高估算的 CO 总量分别为 2.76、2.81、2.94~g，半高处 CO 峰值为 20.0、26.2、163.5~ppm。靠近羽流捕集烟气并提供畅通的补风路径有助于排出 CO。在原始 OP-6 布局、基准烟炱和 0.16~m 垂壁下，排烟量由 0.16 增至 0.64~m$^3$/s 时，远端平均 CO 浓度由 101.04 降至 0.62~ppm。表 1 给出推荐布局的局部暴露指标。

\tabletitle{表 1.\quad 推荐风量 1.12~m$^3$/s 下的半高处暴露指标。各值为细网格八个探头在 120~s 内的极值。}
\begin{reporttable}{L{.18}L{.24}L{.31}L{.27}}
\rowcolor{tablegrey}门态&烟炱产率 (kg/kg)&最低能见度 (m)&最高 CO (ppm)\\
OP-1&0.02&30.00&0.65\\
OP-1&0.05&30.00&0.26\\
OP-1&0.10&30.00&0.28\\
OP-6&0.02&24.15&7.59\\
OP-6&0.05&16.42&11.81\\
OP-6&0.10&4.06&32.83\\
\end{reporttable}
最低能见度和最高 CO 浓度同时出现在 OP-6、烟炱产率 0.10 的 $x=3.6$~m、$t=3.0$~s。最小热分层净高则发生在更远的下游和更晚的时刻，即 OP-6、烟炱产率 0.05 的 $x=6.1$~m、$t=113.0$~s。
